\chapter[Εργαλείο Οπτικοποίησης]{Διαδραστικό Εργαλείο Οπτικοποίησης Αποτελεσμάτων}
\label{ch:dashboard}

Η ταυτόχρονη σύγκριση πολλαπλών μοντέλων και στρατηγικών πρόβλεψης 
σε πολλαπλές χρονικές περιόδους παράγει έναν ιδιαίτερα σύνθετο όγκο δεδομένων. 
Για την αποτελεσματική διαχείριση και ανάλυση των αποτελεσμάτων αναπτύχθηκε μία 
διαδραστική πλατφόρμα οπτικοποίησης (dashboard).
Η πλατφόρμα υλοποιήθηκε σε React με TypeScript και Vite, ενώ τα διαγράμματα αποδίδονται με τη βιβλιοθήκη Recharts.

\section{Συνολική Όψη της Εφαρμογής}
\label{sec:dash_overview}

Η οθόνη οργανώνεται σε τέσσερις ζώνες, που ο χρήστης διαβάζει από πάνω προς τα κάτω
(Σχήμα~\ref{fig:ui_overview}). Στην κορυφή βρίσκεται (1) ο \textbf{πίνακας ελέγχου} με όλα τα
χειριστήρια. Ακολουθεί (2) το \textbf{γράφημα χρονοσειράς} με τις πραγματικές τιμές και τις
προβλέψεις κάθε ορατού μοντέλου σε ωριαία ανάλυση. Κάτω από αυτό βρίσκονται (3) οι
\textbf{πίνακες μετρικών}, που κατατάσσουν χωριστά τα μοντέλα μηχανικής μάθησης και τα μοντέλα
αναφοράς, και τέλος (4) ο \textbf{συγκεντρωτικός πίνακας κατάταξης}, που ενοποιεί όλα τα
μοντέλα σε ενιαία λίστα με την κατηγορία καθενός.

\begin{figure}[H]
  \centering
  \includegraphics[width=1.15\textwidth, center]{ui_screenshots/dash1.png}
  \caption{Συνολική όψη του dashboard με τις τέσσερις ζώνες: (1) πίνακας ελέγχου,
    (2) γράφημα χρονοσειράς, (3) πίνακες μετρικών, (4) συγκεντρωτικός πίνακας κατάταξης.}
  \label{fig:ui_overview}
\end{figure}

\section{Τυπική Ροή Χρήσης}
\label{sec:dash_workflow}

Ο χρήστης διαμορφώνει το πλαίσιο της έρευνας επιλέγοντας το μέγεθος, τη στρατηγική, τη χρονική περίοδο και τα ορατά διαγράμματα στον πίνακα ελέγχου.

\begin{figure}[H]
  \centering
  
\includegraphics[width=\textwidth]{ui_screenshots/panel.png}
  \caption{Ο πίνακας ελέγχου σε μεγέθυνση.}
  \label{fig:ui_panel}
\end{figure}

Οι λειτουργίες συνοψίζονται στον Πίνακα~\ref{tab:dash_controls}.

\begin{table}[H]
  \centering
  \caption{Πίνακας αναφοράς χειριστηρίων του dashboard.}
  \label{tab:dash_controls}
  \small
  \begin{tabular}{p{2.6cm}p{3.4cm}p{7.6cm}}
    \toprule
    \textbf{Χειριστήριο} & \textbf{Επιλογές} & \textbf{Λειτουργία} \\
    \midrule
    Task        & \texttt{Price}, \texttt{Load}                          & Επιλογή τιμής ή φορτίου \\[0.3em]
    Strategy    & \texttt{TF}, \texttt{Rec}, \texttt{MIMO}, \texttt{Direct} & Επιλογή στρατηγικής πρόβλεψης \\[0.3em]
    Period      & \texttt{Dec 2025}, \texttt{Q1 2026}                    & Επιλογή χρονικού ορίζοντα \\[0.3em]
    View        & \texttt{All}, \texttt{ML}, \texttt{Baselines}, \texttt{Best} & Εμφάνιση: \texttt{All} όλα τα μοντέλα, \texttt{ML} μόνο μοντέλα μηχανικής μάθησης, \texttt{Baselines} μόνο μοντέλα αναφοράς, \texttt{Best} καλύτερο μοντέλο και καλύτερο baseline \\[0.3em]
    Models Show & ανά μοντέλο                                            & Εμφάνιση/απόκρυψη καμπυλών \\[0.3em]
    Zoom (From/To) & date picker                                         & Επιλογή συγκεκριμένου χρονικού διαστήματος \\[0.3em]
    $\downarrow$~CSV & ---                                               & Κατέβασμα δεδομένων σε αρχείο csv \\[0.3em]
    Sort by     & \texttt{MAE}, \texttt{RMSE}, \texttt{sMAPE}            & Κατάταξη κατά MAE \\
    \bottomrule
  \end{tabular}
\end{table}

Στην πράξη, σχεδόν όλες οι λειτουργίες φαίνονται στο Σχήμα~\ref{fig:ui_combo}, που
συγκεντρώνει τρεις προβολές της ίδιας εφαρμογής.

\begin{figure}[htbp]
  \centering
  \captionsetup[subfigure]{skip=1pt, font=footnotesize}
  \begin{subfigure}{0.9\textwidth}
    \centering
    \includegraphics[width=\linewidth]{ui_screenshots/combo_top.png}
    \caption{Recursive, Price, Q1 2026, ML, zoom.}
    \label{fig:ui_combo_top}
  \end{subfigure}\\[1pt]
  \begin{subfigure}{0.9\textwidth}
    \centering
    \includegraphics[width=\linewidth]{ui_screenshots/combo_mid.png}
    \caption{MIMO, Price, Baselines, Dec 2025.}
    \label{fig:ui_combo_mid}
  \end{subfigure}\\[1pt]
  \begin{subfigure}{0.9\textwidth}
    \centering
    \includegraphics[width=\linewidth]{ui_screenshots/combo_down.png}
    \caption{Direct, Load, Best, Dec 2025.}
    \label{fig:ui_combo_down}
  \end{subfigure}
  \caption{Τρεις προβολές της εφαρμογής για διαφορετικούς συνδυασμούς εργασίας, στρατηγικής,
    περιόδου και φίλτρου.}
  \label{fig:ui_combo}
\end{figure}

Στην πάνω προβολή η τιμή προβλέπεται με Recursive στο Q1 2026. Το φίλτρο \texttt{ML} κρατά
μόνο τους αλγορίθμους μηχανικής μάθησης και το zoom εστιάζει σε στενότερο διάστημα, ώστε να
φανεί η απόκλιση των καμπυλών από τη συσσώρευση σφάλματος.
Στη μεσαία προβολή η ίδια τιμή
εξετάζεται με MIMO στον Δεκέμβριο 2025 με πατημένο το \texttt{Baselines}, οπότε μένουν μόνο τα
μοντέλα αναφοράς (Naive, Seasonal Profile). 
Στην κάτω προβολή η εργασία γίνεται φορτίο με
Direct και το φίλτρο \texttt{Best} κρατά τρεις καμπύλες, το κορυφαίο μοντέλο, το καλύτερο
baseline και την πραγματική τιμή. 

\begin{figure}[htbp]
  \centering
  \begin{subfigure}{0.52\textwidth}
    \centering
    \includegraphics[width=\textwidth]{ui_screenshots/hover.png}
    \caption{Ωριαία σύγκριση (hover).}
    \label{fig:ui_hover}
  \end{subfigure}\hfill
  \begin{subfigure}{0.42\textwidth}
    \centering
    \includegraphics[width=\textwidth]{ui_screenshots/hmerologio.png}
    \caption{Επιλογή περιόδου.}
    \label{fig:ui_calendar}
  \end{subfigure}
  \caption{Δύο λειτουργίες λεπτομερούς ανάγνωσης: (α) το tooltip του hover για ένα χρονικό
    βήμα και (β) ο επιλογέας ημερομηνίας για μια υποπερίοδο.}
  \label{fig:ui_detail}
\end{figure}


Η χρηστικότητα ενισχύεται από δύο βοηθητικά εργαλεία (Σχήμα~\ref{fig:ui_detail}). Με το hover,
σταματώντας τον δείκτη σε μία ώρα, εμφανίζονται μαζί οι τιμές όλων των ενεργών μοντέλων δίπλα
στην πραγματική, κάτι που καθιστά δυνατό να εξεταστεί ένα spike της τιμής. 
Ο χρήστης έχει τη δυνατότητα να απομονώσει οποιοδήποτε χρονικό διάστημα επιθυμεί για να εστιάσει την ανάλυσή του.

Οι αλλαγές αυτές ενημερώνουν αυτόματα τους δείκτες σφάλματος (MAE, RMSE, sMAPE) στους πίνακες. (Σχήμα~\ref{fig:ui_tables}). Η
κατάταξη χωρίζεται σε μοντέλα μηχανικής μάθησης και μοντέλα αναφοράς, ώστε η σύγκριση να
γίνεται μέσα σε ομοιογενείς ομάδες. Το \texttt{Sort by} ταξινομεί κατά MAE, RMSE ή sMAPE και
το $\downarrow$~\texttt{CSV} κατεβάζει τις τρέχουσες μετρικές σε αρχείο. Στη βάση της σελίδας
βρίσκεται ο συγκεντρωτικός πίνακας με όλα τα μοντέλα της εκάστοτε στρατηγικής.

\begin{figure}[htbp]
  \centering
  \includegraphics[width=0.92\textwidth]{ui_screenshots/pinakes.png}
  \caption{Πίνακες μετρικών: επάνω οι αυτόνομες κατατάξεις των μοντέλων μηχανικής μάθησης
    (αριστερά) και των μοντέλων αναφοράς (δεξιά), και κάτω ο ενιαίος συγκεντρωτικός πίνακας
    όλων των μοντέλων, με επιλογέα ταξινόμησης \texttt{Sort by}.}
  \label{fig:ui_tables}
\end{figure}

Η διάταξη εξυπηρετεί διαφορετικούς χρήστες. Ο ερευνητής εστιάζει σε spikes με το hover και
συγκρίνει στρατηγικές αλλάζοντας tab. Ο αναλυτής της αγοράς ελέγχει μια συγκεκριμένη εβδομάδα
με το zoom και εξάγει τις μετρικές σε CSV. Όποιος θέλει μόνο τη συνολική εικόνα χρησιμοποιεί
το \texttt{Best} και τον ενιαίο πίνακα.

\section{Σύνοψη}
\label{sec:dash_summary}

Στην παρούσα εργασία το dashboard λειτούργησε ως βοηθητικό ερευνητικό εργαλείο. Η άμεση
οπτικοποίηση κάθε συνδυασμού μοντέλου, στρατηγικής και περιόδου επέτρεψε να εντοπιστούν
γρήγορα αδυναμίες, να συγκριθούν στρατηγικές χωρίς εκτέλεση κώδικα και να παρθούν αποφάσεις. 
Στο επόμενο κεφάλαιο, η ποιοτική και οπτική εκτίμηση που εμπλουτίζεται με ποσοτική ανάλυση και στατιστική επαλήθευση.
